% =====================================================================
\chapter{Casos de estudio: edificio, reactor, batería y centro de datos}
\label{tema:5}
% =====================================================================

Este tema reúne los métodos de los anteriores en cuatro problemas de distinta naturaleza: un edificio, donde el objetivo es económico y las restricciones son de confort; un reactor por lotes, donde lo que manda es la seguridad; una batería, donde el problema es casi puramente de optimización, y un centro de datos, el caso industrial publicado más conocido de aprendizaje automático aplicado a la operación de una instalación. Antes se fija cómo comparar controladores de forma honesta. Las fuentes son la revisión de Drgoňa \emph{et al.}~\cite{drgona} y el marco BOPTEST~\cite{boptest} para edificios, la lección~6 de MIT 10.450~\cite{mit10450} y el texto de Seborg \emph{et al.}~\cite{seborg} para el reactor, y las publicaciones de Google y DeepMind~\cite{evansgao,gamblegao} y de Lazic \emph{et al.}~\cite{lazic} para los centros de datos. Los resultados numéricos de los tres primeros casos proceden del cuaderno del módulo.

\section{Cómo comparar controladores}

\begin{definicion}[Indicadores de rendimiento]\label{def:kpi}
Para un periodo de evaluación de $n$ pasos se definen la energía $E=\Dt\sum_kP_k$ (kWh), el coste $\Sigma=\Dt\sum_kc_kP_k$ (€), la violación del confort $D=\Dt\sum_k\bigl(\max\{\underline T_k-T_k,0\}+\max\{T_k-\overline T_k,0\}\bigr)$ (K·h) y, si se quiere un único número, el coste total $\Sigma+\rho D$ con el mismo $\rho$ que usan los controladores.
\end{definicion}

Un protocolo de comparación justo exige: (i) el mismo periodo de prueba, no usado para ajustar ningún controlador; (ii) la misma información, en particular las mismas previsiones: un MPC con la meteorología observada frente a un PID que no la conoce compara información, no métodos; (iii) un controlador de referencia bien ajustado, no un PID descuidado; (iv) los mismos ruidos, y (v) varios periodos o semillas, con una medida de la dispersión. BOPTEST~\cite{boptest} se creó precisamente para normalizar estas comparaciones en edificios: fija los modelos, los escenarios y los indicadores.

\begin{proposicion}[Números aleatorios comunes]\label{prop:crn}
Sean $X$ e $Y$ los costes de dos controladores en un mismo escenario aleatorio, con varianzas $\sigma_X^{2}$, $\sigma_Y^{2}$ y correlación $r$. Entonces $\Var(X-Y)=\sigma_X^{2}+\sigma_Y^{2}-2r\sigma_X\sigma_Y$. Si $\sigma_X=\sigma_Y=\sigma$, evaluar los dos controladores con los mismos escenarios ($r>0$) divide la varianza de la diferencia por $1-r$ respecto a usar escenarios independientes ($r=0$).
\end{proposicion}

\begin{proof}
$\Var(X-Y)=\Var X+\Var Y-2\operatorname{Cov}(X,Y)$ y $\operatorname{Cov}(X,Y)=r\sigma_X\sigma_Y$ (módulo~5, tema~2). Con $\sigma_X=\sigma_Y$, la varianza es $2\sigma^{2}(1-r)$ frente a $2\sigma^{2}$.
\end{proof}

El coste semanal de climatización depende sobre todo del tiempo meteorológico: dos controladores evaluados en la misma semana tienen costes muy correlacionados, y su diferencia se estima con mucha menos incertidumbre que cada coste por separado. El mismo principio guía el entrenamiento del agente del \cref{ej:caso-rl}, que evalúa todos los candidatos de una iteración en los mismos episodios.

\section{Climatización de un edificio}

El caso de la zona ha recorrido los temas~1 a~4. El \cref{tab:hvac} reúne los resultados en la semana de prueba (temperatura exterior media de $3{,}35$ °C) con la planta nominal y con una planta perturbada ($R$ un 20~\% mayor, $C$ un 20~\% menor y COP un 10~\% peor) que ninguno de los controladores conoce.

\begin{table}[htbp]
\centering\small
\begin{tabular}{@{}lcccccc@{}}
\toprule
& \multicolumn{4}{c}{Planta nominal} & \multicolumn{2}{c}{Planta perturbada} \\
\cmidrule(lr){2-5}\cmidrule(l){6-7}
Controlador & $E$ (kWh) & $\Sigma$ (€) & $D$ (K·h) & Ahorro & $\Sigma$ (€) & $D$ (K·h) \\
\midrule
PI & $175{,}4$ & $21{,}06$ & $0{,}00$ & --- & $18{,}75$ & $0{,}00$ \\
MPC, margen $0{,}3$ K & $171{,}1$ & $19{,}28$ & $0{,}00$ & $8{,}4$~\% & $17{,}28$ & $0{,}00$ \\
MPC, sin margen & $166{,}9$ & $18{,}76$ & $3{,}54$ & --- & --- & --- \\
RL (CEM) & $169{,}9$ & $20{,}63$ & $0{,}01$ & $2{,}1$~\% & $18{,}50$ & $0{,}00$ \\
\bottomrule
\end{tabular}
\caption{Climatización de la zona: comparación en la semana de prueba (cuaderno, Partes~2, 4, 6 y~7): energía $E$, coste $\Sigma$ y violación del confort $D$ (\cref{def:kpi}). PI con ajuste SIMC y consigna programada; MPC económico; agente CEM con aleatorización del dominio. El ahorro es en coste respecto al PI.}
\label{tab:hvac}
\end{table}

Los resultados ilustran rasgos generales que la literatura sobre edificios documenta con casos reales~\cite{drgona,wanghong}:
\begin{enumerate}
  \item \emph{El MPC gana porque planifica}: conoce el precio y las ganancias de las próximas 24~h y usa la masa térmica para desplazar el consumo de punta a valle. Su ventaja se mantiene con la planta perturbada (un $7{,}8$~\% menos de coste que el PI), porque realimenta el estado medido en cada paso.
  \item \emph{El margen es parte del diseño}: sin él, el MPC «ahorra» un $2{,}7$~\% más a costa de violar el confort.
  \item \emph{El agente aprende sin modelo explícito, pero con un simulador}: con una política lineal de 8 parámetros mejora al PI y es robusto a la perturbación de la planta gracias a la aleatorización, pero queda lejos del MPC, que dispone de previsiones y de un horizonte de planificación explícito. Una política más rica (una red con previsiones como entrada) y un algoritmo como SAC pueden acortar la distancia, con más datos de simulación y más riesgo de explotar defectos del gemelo.
  \item \emph{El PI sigue siendo la referencia}: no necesita modelo, ni previsiones, ni entrenamiento, y en este caso solo pierde un $8$~\% de coste. En la práctica, el MPC o el agente se colocan \emph{encima} de los PI existentes, enviándoles consignas (\cref{tema:6}).
\end{enumerate}

\begin{table}[htbp]
\centering\small
\begin{tabularx}{\textwidth}{@{}l>{\raggedright\arraybackslash}X>{\raggedright\arraybackslash}X>{\raggedright\arraybackslash}X@{}}
\toprule
& PID & MPC & RL \\
\midrule
Necesita & ajuste local (FOPDT) & modelo identificado y previsiones & simulador fiable (gemelo) y entrenamiento \\
Restricciones & solo límites del actuador & explícitas, en cada paso & penalizadas en la recompensa o con un filtro externo \\
Coste en línea & despreciable & un problema de optimización por paso & una evaluación de la política \\
Adaptación a cambios & reajuste manual & cambiar el modelo o el coste & reentrenar \\
Interpretabilidad & alta & alta (el plan es visible) & baja \\
Cuándo elegirlo & lazos básicos, sin objetivo económico & objetivos económicos con buen modelo y previsiones & dinámicas difíciles de modelar con un simulador rico o decisiones discretas complejas \\
\bottomrule
\end{tabularx}
\caption{Comparación cualitativa de los tres enfoques.}
\label{tab:enfoques}
\end{table}

\section{Control de temperatura de un reactor por lotes}

En un reactor por lotes con una reacción exotérmica, la temperatura sigue una receta (calentar, mantener, enfriar) y el riesgo dominante es el \emph{embalamiento térmico}: si el calor generado por la reacción crece con la temperatura más deprisa de lo que la camisa puede retirarlo, la temperatura se dispara. El análisis de estabilidad de la lección~6 de MIT 10.450~\cite{mit10450}, que compara las curvas de calor generado y calor retirado, se resume en el criterio siguiente.

\begin{definicion}[Modelo del reactor]\label{def:reactor}
Para una reacción de primer orden $\mathrm A\to\mathrm B$ en un reactor de volumen $V$ refrigerado por una camisa a temperatura $T_j$:
\[
\begin{gathered}
\dot C_A=-k(T)C_A,\qquad \rho c_pV\,\dot T=Q_g(T,C_A)-Q_r(T),\\
Q_g=(-\Delta H)\,k(T)\,C_AV,\qquad Q_r=UA\,(T-T_j),
\end{gathered}
\]
con $k(T)=k_{\mathrm{ref}}\exp\bigl(-\frac{E}{R}(\frac1T-\frac1{T_{\mathrm{ref}}})\bigr)$ (ley de Arrhenius).
\end{definicion}

\begin{proposicion}[Criterio de Semenov]\label{prop:semenov}
Sea $C_A$ aproximadamente constante (fase inicial del lote, con conversión pequeña) y sea $\bar T$ un equilibrio de la ecuación de energía, $Q_g(\bar T)=UA(\bar T-T_j)$. El equilibrio es asintóticamente estable si $Q_g'(\bar T)<UA$ e inestable si $Q_g'(\bar T)>UA$. Además, $Q_g'(T)=Q_g(T)\,E/(RT^{2})$.
\end{proposicion}

\begin{proof}
Con $C_A$ fijo, la temperatura cumple la ecuación escalar $\dot T=F(T)=(Q_g(T)-UA(T-T_j))/(\rho c_pV)$. Por la \cref{prop:lineal}, el linealizado en $\bar T$ es $\dot{\delta T}=\lambda\,\delta T$ con $\lambda=F'(\bar T)=(Q_g'(\bar T)-UA)/(\rho c_pV)$; si $\lambda<0$ las desviaciones decaen y si $\lambda>0$ crecen (teorema de linealización de Lyapunov~\cite[cap.~5]{am}). La derivada de la exponencial da $k'(T)=k(T)\,E/(RT^{2})$, de donde se sigue la última fórmula.
\end{proof}

\begin{ejemplo}[Caso de estudio: margen térmico de un reactor]\label{ej:caso-reactor}
\emph{Datos.} $V=1$ m³, $\rho c_pV=4000$ kJ/K, $-\Delta H=100$ kJ/mol, $C_{A,0}=2000$ mol/m³, $E/R=10\,000$ K, $k_{\mathrm{ref}}=10^{-4}$ s$^{-1}$ a $T_{\mathrm{ref}}=350$ K y $UA=2$ kW/K.

\emph{Punto de operación.} A $350$ K y conversión nula, $Q_g=100\cdot10^{-4}\cdot2000=20$ kW; el equilibrio exige $T_j=350-20/2=340$ K. La pendiente es $Q_g'=20\cdot10^{4}/350^{2}=1{,}633$ kW/K, menor que $UA=2$ kW/K: el punto es estable, con constante de tiempo $4000/(2-1{,}633)=10\,889$ s, unas 3~h.

\emph{Margen.} Con $T_j=340$ K hay un segundo equilibrio en $T_u=355{,}23$ K, donde $Q_g'=2{,}414>UA$: es inestable y separa las condiciones iniciales que vuelven a $350$ K de las que se embalan. En el cuaderno (Parte~9), partiendo de $354$ K la temperatura vuelve hacia $350$ K, y partiendo de $356$ K supera los $400$ K en unas $4{,}6$ h. Además, existe una temperatura de camisa crítica: las curvas $Q_g$ y $Q_r$ son tangentes cuando $Q_g(T)=UA\,T^{2}R/E$, es decir, en $T^{*}=352{,}69$ K con $T_j^{*}=T^{*}-RT^{*2}/E=340{,}26$ K; por encima de ella no existe ningún equilibrio y la fase inicial se embala desde cualquier temperatura. El margen de operación es de solo $0{,}26$ K de camisa y $5{,}2$ K de reactor.

\emph{Elección del método de control.} El esquema clásico es una \emph{cascada}: un PID maestro sobre la temperatura del reactor fija la consigna de un PID esclavo sobre la camisa, que rechaza rápido las perturbaciones del agua de refrigeración~\cite{mit10450,seborg}. Un NMPC con el modelo de la \cref{def:reactor} puede seguir la receta minimizando el tiempo de lote con restricciones explícitas $T\le T_{\max}$ y $T_j\le T_j^{*}-\delta$. Para un agente de RL, cualquier exploración en planta es inaceptable; se entrena en el gemelo y se despliega detrás de un \emph{filtro de seguridad} que limita $T_j$ y fuerza refrigeración máxima si $T$ se acerca a $T_u$ (cuaderno, Parte~9; \cref{tema:6}).
\end{ejemplo}

\section{Gestión de una batería con precios horarios}

El \cref{ej:bateria-lp} y el \cref{ej:bateria-mdp} resuelven el mismo problema por dos vías. Con precios conocidos, el LP es exacto y rápido: $1{,}8842$ € en el día del cuaderno. La programación dinámica sobre una rejilla de 81 niveles obtiene $1{,}8837$ €, y su interés aparece cuando los precios son \emph{inciertos}: la PD estocástica maximiza el beneficio esperado con el precio como parte del estado, y el RL la aproxima cuando el estado crece (varias baterías, una planta fotovoltaica, la degradación como estado). Tres decisiones de modelado dominan el resultado: el rendimiento de ida y vuelta (un ciclo solo compensa si el cociente de precios supera $1/\eta^{2}$), el coste de degradación por kWh cíclico y la condición terminal (sin $e_{24}=e_0$, el optimizador vacía la batería al final del día porque el modelo no ve el día siguiente).

\section{Refrigeración de centros de datos: el caso publicado de Google y DeepMind}

Lo que sigue usa exclusivamente la información publicada; los detalles internos del sistema no se han hecho públicos.

\emph{2016: recomendaciones.} Según la publicación de DeepMind de julio de 2016~\cite{evansgao}, un conjunto de redes neuronales profundas se entrenó con datos históricos de miles de sensores (temperaturas, potencias, velocidades de bombas, consignas) para predecir la eficacia en el uso de la energía (PUE, cociente entre la energía total del edificio y la consumida por los equipos informáticos) futura; otros dos conjuntos de redes predecían la temperatura y la presión de la hora siguiente para comprobar que las acciones recomendadas respetaban los límites de operación. En pruebas en un centro de datos en funcionamiento, alternando periodos con el sistema activado y desactivado, se publicó una reducción de hasta el 40~\% de la energía de refrigeración, equivalente a un 15~\% de reducción del exceso de PUE sobre la unidad.

\emph{2018: control autónomo.} En agosto de 2018~\cite{gamblegao} se publicó que el sistema pasó a controlar directamente la refrigeración: cada cinco minutos, un sistema de IA en la nube toma una instantánea de miles de sensores, predice con redes neuronales el efecto de combinaciones de acciones sobre el consumo futuro y elige las que lo minimizan respetando restricciones de seguridad; el sistema de control local verifica las acciones antes de aplicarlas. Entre las salvaguardas publicadas figuran la estimación de la incertidumbre (solo se aplican acciones de alta confianza), la verificación en dos capas, la vuelta automática a un estado neutro si se violan restricciones, reglas heurísticas de respaldo y la posibilidad de que los operadores salgan del modo de IA en cualquier momento. Se publicó un ahorro medio de alrededor del 30~\% tras unos meses de operación.

\emph{Un estudio académico relacionado.} Lazic \emph{et al.}~\cite{lazic} describieron en NeurIPS 2018 un MPC con un modelo ARX disperso de historia de 5 pasos, ajustado por mínimos cuadrados recursivos tras unas 3~h de exploración aleatoria acotada, que regulaba temperaturas y presiones de un sector de servidores cada 30~s; la exploración limitaba el rango y la variación máxima de cada actuador entre pasos. En condiciones comparables, el coste del controlador fue de $116{,}7$ frente a $142{,}2$ de los controladores locales, y estimaron un ahorro de alrededor del 9~\% en el coste de refrigeración de ese sector, con la advertencia de que harían falta experimentos más largos para confirmarlo.

El caso reúne casi todo el módulo: un modelo aprendido (\cref{tema:2}), un optimizador sobre predicciones con restricciones (\cref{tema:3}), un periodo de recomendación con humanos en el lazo antes del control autónomo, salvaguardas en capas y la medición del ahorro con experimentos de activación alterna (\cref{tema:6}).

\begin{table}[htbp]
\centering\footnotesize
\begin{tabularx}{\textwidth}{@{}>{\raggedright\arraybackslash}p{3.1cm}>{\raggedright\arraybackslash}X>{\raggedright\arraybackslash}p{2.0cm}>{\raggedright\arraybackslash}X@{}}
\toprule
Herramienta & Para qué sirve & Licencia o tipo & Cuándo elegirla \\
\midrule
EnergyPlus~\cite{crawley} & simulación energética de edificios completa & libre & gemelo físico de un edificio \\
Biblioteca Buildings de Modelica~\cite{modelica} & componentes de climatización y control & libre & modelos de instalaciones a medida \\
BOPTEST~\cite{boptest} & bancos de prueba normalizados para controladores de edificios & libre & comparar controladores de forma reproducible \\
Sinergym~\cite{sinergym} & EnergyPlus como entorno Gymnasium & libre (MIT) & entrenar agentes de climatización \\
CasADi, do-mpc~\cite{casadi,dompc} & NMPC del reactor con restricciones de seguridad & libres (LGPL-3) & procesos químicos no lineales \\
CVXPY, Pyomo & programación lineal de baterías y despacho & libres & arbitraje y planificación energética \\
\bottomrule
\end{tabularx}
\caption{Tecnologías para los casos de estudio.}
\label{tab:tec5}
\end{table}

\begin{observacion}[Errores frecuentes en la práctica]\label{obs:errores5}
(i) \emph{Referencia débil}: comparar con un PID mal ajustado infla cualquier ahorro. (ii) \emph{Información desigual}: dar previsiones perfectas a uno de los controladores. (iii) \emph{Una sola semana o una sola semilla}, sin dispersión. (iv) \emph{Confundir porcentajes}: un 40~\% de la energía de refrigeración no es un 40~\% de la factura (\cref{eje:5-pue}). (v) \emph{Ignorar el coste de implantación}: sensores, modelado, mantenimiento del modelo y supervisión cuestan más que el algoritmo. (vi) \emph{En reactores, optimizar el rendimiento sin restricciones de seguridad explícitas}, o confiar a la penalización de la recompensa lo que debe garantizar un enclavamiento.
\end{observacion}

\begin{lecturas}
\item J.~Drgoňa \emph{et al.}~\cite{drgona} y Z.~Wang y T.~Hong~\cite{wanghong}, sobre MPC y RL en edificios.
\item D.~Blum \emph{et al.}~\cite{boptest}, «Building optimization testing framework (BOPTEST)».
\item MIT 10.450~\cite{mit10450}, lección~6 «The Exothermic Reactor».
\item R.~Evans y J.~Gao~\cite{evansgao}; C.~Gamble y J.~Gao~\cite{gamblegao}; N.~Lazic \emph{et al.}~\cite{lazic}.
\end{lecturas}

\section*{Ejercicios del tema 5}
\addcontentsline{toc}{section}{Ejercicios del tema 5}

\begin{ejercicio}{1}{eje:5-kpi}
Un proveedor anuncia que su controlador «ahorra un 25~\% de energía» en un edificio. Enumere las preguntas que haría antes de creerlo, relacionándolas con el protocolo de comparación de este tema.
\end{ejercicio}

\begin{ejercicio}{2}{eje:5-ahorro}
Con el \cref{tab:hvac}: (a) calcule el ahorro de energía y de coste del MPC y del agente respecto al PI; (b) calcule el coste total con $\rho=2$ €/(K·h) del MPC sin margen y diga a partir de qué valor de $\rho$ sería preferible al MPC con margen; (c) explique por qué el MPC consume casi la misma energía que el PI pero cuesta menos.
\end{ejercicio}

\begin{ejercicio}{2}{eje:5-crn}
El coste semanal de un controlador tiene desviación típica de $2$ € entre semanas. (a) Calcule la desviación típica de la diferencia de costes entre dos controladores evaluados en semanas independientes y en las mismas semanas con correlación $r=0{,}9$. (b) ¿Cuántas semanas independientes harían falta para igualar la precisión de una sola semana común?
\fuente{adaptado de D.~C. Montgomery, \emph{Design and Analysis of Experiments}, cap.~2 «Simple Comparative Experiments», ventajas del diseño de comparaciones pareadas~\cite{montgomery}.}
\end{ejercicio}

\begin{ejercicio}{2}{eje:5-semenov}
Con los datos del \cref{ej:caso-reactor}: (a) calcule $Q_g$ y $Q_g'$ a $360$ K con $C_A=C_{A,0}$; ¿qué valor de $UA$ haría falta para que un equilibrio a esa temperatura fuera estable? (b) Por ensuciamiento, $UA$ baja a $1{,}5$ kW/K. ¿Qué $T_j$ mantiene el equilibrio a $350$ K y es estable? (c) ¿Qué medida de operación propondría?
\fuente{adaptado de MIT 10.450, primavera de 2006, lección~6 «The Exothermic Reactor», §6.8 (calor generado frente a calor retirado)~\cite{mit10450}.}
\end{ejercicio}

\begin{ejercicio}{2}{eje:5-pue}
Según la publicación de 2016~\cite{evansgao}, reducir un 40~\% la energía de refrigeración equivalió a reducir un 15~\% el exceso de PUE. (a) ¿Qué fracción del exceso de PUE representaba la refrigeración? (b) Con un PUE ilustrativo de $1{,}12$, ¿cuál sería el PUE resultante y la reducción porcentual del consumo total del centro? (c) Comente la diferencia entre los tres porcentajes.
\end{ejercicio}

\begin{ejercicio}{2}{eje:5-bateria}
La batería del \cref{ej:bateria-lp} ($1$ kWh, $1$ kW, $\eta=1$) se opera con precios previstos $(0{,}10;\ 0{,}30;\ 0{,}05;\ 0{,}25)$, pero los precios reales resultan $(0{,}10;\ 0{,}30;\ 0{,}28;\ 0{,}25)$. (a) ¿Qué beneficio obtiene el plan calculado con la previsión si se aplica en lazo abierto? (b) ¿Y un MPC que en cada hora conoce el precio de esa hora y reoptimiza con la previsión para las restantes? (c) ¿Cuál habría sido el óptimo con información perfecta?
\end{ejercicio}

\begin{ejercicio}{3}{eje:5-semenov-dem}
(a) Demuestre que la condición de tangencia entre $Q_g$ y $Q_r$ conduce a $T_j^{*}=T^{*}-RT^{*2}/E$, con $T^{*}$ solución de $Q_g(T)=UA\,RT^{2}/E$. (b) Demuestre que, si existen exactamente dos equilibrios y $Q_g$ es convexa en el intervalo que los contiene, el inferior es estable y el superior inestable. (c) Explique por qué la aproximación de $C_A$ constante es conservadora para la fase inicial del lote.
\fuente{adaptado de MIT 10.450, primavera de 2006, lección~6 «The Exothermic Reactor», §6.8 y §6.10 (estados estacionarios múltiples y su estabilidad)~\cite{mit10450}.}
\end{ejercicio}

\begin{ejercicio}{3}{eje:5-protocolo}
Diseñe el protocolo para comparar en un edificio de oficinas real el PID existente con un MPC durante un invierno: cómo asignar los días a cada controlador, qué indicadores medir y con qué normalización (grados-día, ocupación), qué análisis estadístico aplicar y qué criterios de parada por seguridad o confort establecer.
\end{ejercicio}
